% Cote 119: offprints, reports, courses, CV and letters (typed and printed part): archive cote 119, site page 186 (PDF page 187).
% typescript. Draft compared: none.
% AI-transcribed by Claude Opus 5.5 (claude-opus-5-5, adaptive thinking, effort high), Message Batches API, batch r119_full (msgbatch_01DwZNm5Y7YsKBWyAN4sFs8w), 2026-11-13; not yet checked by eye.
\scanpage{119/186}
\mspage{- 5 -}

un point $x_o$ tel qu'il existe un entier $( \geq 1$ $k$ et un voisinage de $x_o$ sur
lequel $f(x) = \dfrac{f_1(x)}{(x-x_o)^k}$ avec $f_1(x_o) \neq 0$ et $f_1$ analytique$]$ .

8.1.2. Au voisinage de tout point $\hand{U}$ $f$ peut se mettre sous la forme du quo-%
tient de deux fonctions analytiques $\dfrac{k(x)}{g(x)}$ , le dénominateur n'étant pas iden-%
tiquement nul.

8.1.3. Si $f$ est \unsure{méromorphe} dans $\hand{U}$ et non identiquement nulle $\frac{1}{f}$ est méromor-%
phe dans $\hand{U}$ . En effet d'après le principe du prolongement analytique $f$ ne
peut être nulle au voisinage d'aucun point de $\hand{U}$. Donc en tout point de $\hand{U}$ $\dfrac{1}{f(x)}$
est analytique ou possède un pôle en ce point. On a donc :

8.1.4. \uline{Proposition}.

L'ensemble des fonctions méromorphes sur un ouvert connexe $\hand{U}$ du
plan complexe forment un corps.

Il suffit de rajouter 8.1.3. à la proposition 1.3. du fascicule.
Fonctions Analytiques.

8.2. \uline{Séries de Laurent}.

8.2.1. \uline{Définition}.On appelle série de Laurent une série de la forme $\sum\limits_{n \hand{\in \mathbb{Z}}} a_n z^n$ .

C'est donc la somme d'une série entière $\sum\limits_{n \hand{\in \mathbb{N}}} a_n z^n$ de rayon de convergence $\rho_1$
et d'une série de la forme $\sum\limits_{n<0} a_n z^n$. Soit $(\rho_2)^{-1}$ le rayon de convergence de
la série entière $\sum\limits_{n<0} a_n z^{-n}$.

On a les séries convergentes :
\[ f_1(z) = \sum_{n \geq 0} a_n z^n \quad \text{pour} \quad |z| < \rho_1 \]
\[ f_2\ (z) = \sum_{n < 0} a_n z^n \quad \text{pour} \quad |z| > \rho_2 \]

$f_1(z)$ est évidemment holomorphe. Montrons que $f_2(z)$ est également
holomorphe. Posons $z = \frac{1}{u}$ .

\margin{186}
